%============================ 3. SETUP AND CALIBRATION ========================
\isec{\faTools}{Experimental Setup and Calibration}
\label{sec:setup}

\isub{\faCamera}{Camera pose}
The range calibration is only valid at one arm pose, so the arm is commanded to the saved pose
$(x,z)=(0.123,0.067)$\,m before every run and the controller refuses to drive if it is more than 6\,mm off.
Because the tag lies flat, its normal is the floor normal, and the detected tag pose gives the camera height and pitch
directly: \textbf{0.315\,m and 6.1$^\circ$ down} (Phase 2 had 5.3$^\circ$; the arm had been moved since). The pitch
$\beta=6.1^\circ$ is used in the level-frame rotation of Section~\ref{sec:method}.

\isub{\faRulerCombined}{Range calibration with the robot driving itself}
The camera's forward offset $f$ is mapped to the tape distance (chassis front to tag centre) by a straight line
$d = s f + o$. Instead of placing the robot by hand at five distances, a script
(\code{scripts/auto_calibrate_range.py}) drives the robot at 5\,cm/s, steering at the tag, and stops at
1.20, 1.00, 0.80, 0.60 and 0.50\,m; at each stop it averages 5\,s of camera readings. Only the first and last
positions need the tape; the stops in between are known from odometry and serve as a check.

\begin{table}[H]
\centering\small
\rowcolors{2}{kfupmlight}{white}
\begin{tabular}{@{}lrrrl@{}}
\rowcolor{kfupm}\hd{Stop} & \hd{Tape / odometry (m)} & \hd{Camera $f$ (m)} & \hd{Spread (mm)} & \hd{Use}\\
start & 1.200 (tape) & 1.3219 & 1.4 & fit anchor\\
1 & 0.9915 (odometry) & 1.1190 & 1.3 & check: fit residual $+3.2$\,mm\\
2 & 0.7787 (odometry) & 0.9120 & 0.6 & check: fit residual $+6.6$\,mm\\
last & 0.470 (tape) & 0.6003 & 0.2 & fit anchor\\
\end{tabular}
\caption{Range calibration, run 2 (\code{results/calibration/}). The fit through the two tape anchors gives
$s=1.0115$, $o=-0.1372$\,m for a 155\,mm tag size; with the measured 154\,mm the camera offsets scale by 154/155, so
$s=1.0181$ is used with the same $o$.}
\label{tab:calib}
\end{table}

\begin{itemize}
  \item \textbf{Run 1 was discarded.} Without steering the robot drifted on the rug and ended with the tag
        18$^\circ$ to the side; the camera gives the forward offset while the tape measures the diagonal, about
        2\,cm different at that angle. Steering was added to the script for run 2 (bearing stayed within
        0.2 to 0.4$^\circ$).
  \item \textbf{Run 2 had a Wi-Fi drop} of 3\,s during the leg to 0.60\,m; the robot kept its last speed and overshot
        to 0.47\,m, and \code{/odom} restarted from zero, so odometry after the drop was not used. This is why the
        fit uses the two tape anchors only.
  \item The two odometry check points sit 3 and 7\,mm above the fit, i.e.\ odometry over-reads the distance travelled
        by about 1.5\,\% on this rug (run 1 gave 2.1\,\%). This is the along-track drift the camera has to correct.
\end{itemize}

\isub{\faWrench}{Odometry: two driver problems found and fixed}
\label{sec:odomfix}
Three test drives (trial 0, kept in \code{results/test_runs/}) were made before the 10 trials.

\begin{table}[H]
\centering\small
\setlength{\tabcolsep}{4pt}
\rowcolors{2}{kfupmlight}{white}
\begin{tabularx}{\linewidth}{@{}l>{\raggedright\arraybackslash}Xrrrrr@{}}
\rowcolor{kfupm}\hd{Test} & \hd{Odometry source} & \hd{$d_{est}$} & \hd{camera} & \hd{wheels only} & \hd{tape} & \hd{$e$ (mm)}\\
1 & \code{/odom} twist ($v,\omega$ fields) & 0.514 & 0.535 & 0.158 & 0.540 & $-26$\\
2 & speed from \code{/odom} pose; one glitch at the stop & 0.614 & 0.500 & 0.507 & 0.520 & $+94$\\
3 & pose + plausibility guard (final setting) & 0.497 & 0.497 & 0.520 & 0.490 & $+7$\\
\end{tabularx}
\caption{Test drives that shaped the odometry input (distances in m).}
\label{tab:testruns}
\end{table}

\begin{itemize}
  \item \textbf{Twist is wrong and late.} At a true speed of 0.15\,m/s (from the camera) the twist field read
        0.18 to 0.21\,m/s, and at 0.03\,m/s it read 3 to 5 times too much, lagging by more than a second. Integrated, it
        put the robot 38\,cm past its real position. The \code{/odom} \emph{pose}, in contrast, agreed with the camera to
        1 to 2\,\% in both calibration runs, so the speed and turn rate are now differentiated from consecutive poses
        (\code{odom_source: pose}).
  \item \textbf{Pose glitches.} While standing still, one pose step gave $-0.72$\,m/s. The EKF moved 11\,cm, and the
        $\chi^2$ gate then rejected every correct camera reading (NIS 12.5) because they disagreed with the corrupted
        estimate. A plausibility guard now skips any odometry step faster than 0.40\,m/s or 1.5\,rad/s (more than
        twice what this controller ever commands); skipped steps are counted in \code{n_odom_glitches}.
\end{itemize}

\isub{\faClipboardList}{Protocol of a trial}
\begin{enumerate}
  \item Place the robot on the taped start marks (front edge on the 1.20\,m line, nose at the tag); arm at the saved pose.
  \item Run \code{ros2 launch rm_odom_ekf phase3.launch.py trial:=N}. Do not touch the robot until \stbadge{DONE}.
  \item Without moving the robot, tape from the middle of the chassis front to the tag centre; record with
        \code{scripts/add_measurement.py N} \textit{reading}.
  \item A trial is invalid if a safety stop ends it; it is then repeated with the same $N$.
\end{enumerate}
The configuration (\code{config/robot.yaml}) was not changed during the 10 trials.
